\section{Discussion}
\label{sec:discussion}

\paragraph{What the platform enables.}
The comparisons this research area needs are now configuration rather than reimplementation: contrasting intervention presentation forms, commit boundaries, or decision strategies under identical sensing is a matter of session templates and module selection, with counterbalancing supplied by the platform's randomised item ordering and evidence supplied by the session record.
The unbundled study that our own confound calls for, crossing the revised restore with and without the highlight cue, is exactly such a configuration.
Beyond single studies, the exported records form a corpus on which an external decision provider can be trained offline and then attached through the same seam it will be judged at, closing the loop the programme's envisaged AI-driven adaptation requires \cite{dtucompute2025rtr}.

\paragraph{Measure the quality attribute you claim.}
The over-repositioning result is the strongest argument for a design principle the platform embodies: a system should measure the property it claims to provide, so it can report its failures as readily as its successes.
Context preservation was a headline feature, yet its original geometry made small reflows worse, and nothing in ordinary operation would have surfaced that; the per-intervention residual in the session record did.
The same record also surfaced the repeated derived-event payloads and the mismatch between the two sweep runs, which a success-only log would have hidden.
We suggest the principle transfers beyond reading platforms: any interface that changes state during a continuous user task can hold the change until a natural boundary, re-anchor the user's focus across it, and record what the change cost.

\paragraph{Limitations.}
The behavioural comparison rests on four participants drawn from the authors and acquaintances, reading two texts in an unbalanced order under researcher-timed manual interventions; the numbers are descriptive, and because the \on{} condition bundles the restore with the highlight cue, no causal attribution to the restore geometry is possible.
The revised restore is verified geometrically, in a sweep whose two runs are not matched, and has not been re-evaluated with readers.
The provider contract has one external integration, written by us around a collaborator's code, and the automated decision path was exercised in a single advisory run on one host; the runtime neither falls back to the built-in strategy on provider detachment nor enforces heartbeat expiry.
Timing figures are loopback figures that stop short of the display, the exports do not pin the full browser, font, and provider configuration, and our participants were general readers, so no claim here extends to the vision-impaired or dyslexic readers who motivate the programme.

\paragraph{Future work.}
Three studies follow directly.
First, the three-condition study separating restore geometry from highlighting on the revised mechanism, which would also replace the authors-as-participants sample and should pin texts, order, and intervention timing.
Second, a campaign-scale run of the automated decision path with a learned provider, turning the single-run check into a measured result, ideally with an end-to-end latency measurement to the display.
Third, the properly powered efficacy study the platform was built to host, asking whether specific interventions improve reading rather than only measuring their immediate cost.
On the platform side, adding automatic fallback and heartbeat enforcement with a fault-injection test and attracting external providers and multimodal sources behind the existing ports would push the contract from author-validated toward independently validated.
